% PAPER_FERROMAGNETO --- standalone Letter targeting Physical Review Letters.
%
% The "form/scale" Letter. Central claim (version 2, corrected by the Phase-1
% evaluation in PAPER_FERROMAGNETO_GAPS.md):
%   A single Poisson causal substrate derives DIMENSIONLESS structure (forms,
%   ratios, exponents) and INTERNAL scales in lattice units, across four domains
%   that were not built to agree (geometry, matter, cosmology, quantum collapse);
%   what is never fixed by the substrate is the lattice->SI conversion, so every
%   absolute physical scale (G, hbar, a0, f_pi, eta) needs one external calibration.
%
% Honest framing enforced throughout (see the charter's writing rules):
%   - NOT "cannot generate scale": the substrate DOES generate internal scales
%     (h_sat, M_Skyrmion, sigma, lambda_Sk in lattice units); the lattice->SI map
%     is the single external input.
%   - "causal substrate" is the unifying object; the ferromagnet is the sector that
%     carries matter+cosmology; the collapse domain uses the BARE causal spectrum.
%   - cosmology link (nu = -0.37, not -0.5) declared [WEAK], not hidden.
%   - collapse domain = UNPLANNED CONVERGENCE with Stochastic Rupture, not derivation.
%   - QCD analogy includes the honest difference: QCD transmutes Lambda_QCD from a
%     dimensionless coupling; in TEIC no such transmutation mechanism is identified.
%   - guard caveat declared: the anti-circularity guard forbids inserting scales, so
%     the external-ness is partly methodological.
%
% Compiles with pdflatex+bibtex (revtex4-2 + apsrev4-2.bst) against the shared
% TEIC_CONSOLIDATED.bib (with added ZambuziSR2026 + QCD-transmutation entries).
%
% Reframe (jul/2026, PLANO_REENQUADRAMENTO): (1) the outdated "a non-equilibrium
% route is not ruled out" corrected -- the NESS parametric drive was tested and is
% mean-field too (amplitude sweep killed the apparent 0.53 exponent); (2) the scale
% closure now carries its structural reason (Campbell-Mecke / boost non-compactness
% for invariant rules; tree/Bethe for the non-invariant escapes); (3) legacy
% citations PaperI/II/V retargeted to the active companion series
% (MatterGravityPRD/SU3PRD); PaperIV kept (sole record of the chi_par MOND-form
% measurement). No positive result changed.
\documentclass[aps,prl,reprint,nofootinbib,superscriptaddress,floatfix]{revtex4-2}
\usepackage{amsmath,amssymb}
\usepackage{graphicx}
\usepackage{bm}
\usepackage{booktabs}
\usepackage{hyperref}

\begin{document}

\title{Dimensionless structure across four domains of a causal network:\\
forms derive, absolute scales do not}

\author{Miqueias Alves Mendes}
\email{mikalvesm@gmail.com}
\affiliation{Independent Researcher, Ibiapina, Cear\'a, Brazil}
\date{\today}

\begin{abstract}
A single discrete causal substrate---a Poisson sprinkling of events with their causal
order---reproduces, under a guard that forbids inserting any relativistic or
dimensionful expression into the generating code, the \emph{dimensionless} structure of
physics in four domains that were not built to agree: emergent Newtonian gravity, a
spin-$\tfrac12$ baryon with $SU(3)$ colour confinement, the deep-MOND interpolation, and
the order parameter of objective wave-function collapse. In every domain the substrate
also fixes \emph{internal} scales in lattice units (the gravitational response, the
soliton mass and string tension, the MOND response knee), but it never fixes the
conversion of those lattice units to physical (SI) units, so $G$, $\hbar$, $a_0$, $f_\pi$
and the collapse threshold $\eta$ each require one external calibration. The pattern is
the lattice-QCD situation, with one honest difference: lattice QCD generates its scale
$\Lambda_{\rm QCD}$ from a dimensionless coupling by dimensional transmutation, whereas
here no transmutation mechanism operates: we tested the cheapest equilibrium signature---a
divergent correlation length at the ordering transition---and the substrate has none (it is
mean-field, in equilibrium and under a non-equilibrium parametric drive), for a structural
reason: for any Lorentz-invariant connection rule on the sprinkling the expected
coordination diverges with system size, so no critical length survives to anchor a
manufactured scale, and the internal scales ride on a fixed action normalisation. The claim is falsifiable: a single domain in which the
lattice$\to$SI conversion is fixed by the substrate, with no external input, would refute
it. We report the supporting evidence and state plainly where it is weak (the cosmology
link) and where it is convergence rather than derivation (the collapse domain).
\end{abstract}

\maketitle

\emph{The claim.}---Emergent-spacetime programmes are often read as promising to derive
``everything.'' We make and test a deliberately narrower, falsifiable statement. A single
Poisson causal substrate \cite{BLMS1987,BombelliHensonSorkin2009,Surya2019}, equipped
with one internal field and analysed under an anti-circularity guard, derives the
\emph{dimensionless} content of physics---functional forms, exponents, ratios, and
topological facts---in four domains that share no design and were measured in separate
campaigns. The same substrate fixes internal scales, but only in lattice units; the one
thing it does not supply is the conversion of those units to physical ones. The result is
a sharp boundary, identical in all four domains: \emph{forms and dimensionless ratios are
derived; the lattice$\to$SI map---and hence every absolute scale $G$, $\hbar$, $a_0$,
$f_\pi$, $\eta$---is external.} This is interesting precisely because it is not a ``theory
of everything'' statement: it is specific, it is falsifiable (below), and it recurs in
domains, such as quantum-measurement collapse, that were never intended to share a vacuum
with galactic dynamics.

\emph{The substrate and the ferromagnet.}---The substrate is a causal set: a locally
finite set of events with a transitive, irreflexive partial order $\prec$, realised by
Poisson sprinkling into $(3{+}1)$-dimensional Minkowski space \cite{BLMS1987}, breaking no
Lorentz frame on average \cite{BombelliHensonSorkin2009}. On three of the four domains the
relevant object is the \emph{causal ferromagnet}: an internal orientation $\vec n$ (in
$O(3)$, $SU(2)$, or $SU(3)$) on each event, coupled across causal links by the minimal
tension $S=J\sum_{\langle ij\rangle}\Delta\tau_{ij}[1-\vec n_i\!\cdot\!\vec n_j]$, which
orders spontaneously through a continuous transition (genuine by finite-size scaling,
Binder cumulant $U_4=2/3$) \cite{GoldstonePRD}. The fourth domain (collapse) uses the
\emph{bare} causal spectrum, with no ferromagnet---so the unifying object is the causal
substrate, of which the ferromagnet is one sector. Throughout, a guard tokenises the
generating code and fails if a Lorentz factor, a metric redshift, a value of $c$, $G$, or
any dimensionful constant appears outside a labelled, non-generating comparison block.
This guard has a consequence we state up front and return to: by forbidding the insertion
of scales, it guarantees that only dimensionless structure \emph{can} emerge---so part of
the ``scales stay external'' pattern is a property of the discipline, not only of the
substrate.

\emph{Four domains, one pattern.}---Table~\ref{tab:pattern} is the central result: for
each domain, what the substrate derives (form), what internal scale it fixes in lattice
units, and which physical scale remains external. We summarise the evidence; the technical
detail lives in the cited domain papers.

\begin{table*}[t]
\caption{One pattern in four domains of the causal substrate. ``Form derived'' lists the
dimensionless structure that emerges without calibration; ``internal scale'' is a
substrate-fixed quantity in \emph{lattice} units; ``external (SI)'' is the absolute
physical scale that requires one calibration. We distinguish two kinds of external scale: a
dimensional SI constant with no computational route ($G$, $f_\pi$, $\hbar$), and a
\emph{geometric} one ($\dagger$)---a calculable graph property whose form and mechanism are
derived and whose only missing input is a finite-size/clustering correction, not an
arbitrary world constant (here the collapse threshold $\eta$, a percolation threshold fixed
by the causal graph's degree heterogeneity and clustering). The cosmology form is weak
(declared); the collapse domain is unplanned convergence with an independent collapse
theory, not a derivation.}
\label{tab:pattern}
\resizebox{\textwidth}{!}{%
\begin{tabular}{lllll}
\toprule
Domain & Form derived (dimensionless) & Internal scale (lattice units) & External (SI) & Status \\
\midrule
Geometry \cite{MatterGravityPRD,TEICDoc1}
 & $\theta\propto1/r$ (exp.\ $-0.992$); Poisson; superpos.\ $10^{-9}$; $\sqrt{1{-}2M/r}$ to $0.2\%$
 & $G_{\rm net}=A/M=0.93$, $\propto1/K$
 & $G$ (via $K$, granularity)
 & solid \\
Matter \cite{MatterGravityPRD,SU3PRD,GoldstonePRD}
 & $B{=}1$, spin-$\tfrac12$; $c{=}\sqrt{2/3}$; $\mu_p/\mu_n{=}{-}1.51$; $[4,16,36]$; $V{\sim}\sigma r$
 & $M_{\rm Sk}{\approx}146{-}207$, $\sigma$, $\lambda_{\rm Sk}{=}a/\sqrt{120}$
 & $f_\pi$ (GeV); $e$ one calibration
 & solid \\
Cosmology \cite{PaperIV}
 & deep-MOND $\chi_\parallel\!\sim\!h^{-0.37}$ ($-0.4{\pm}0.1$; ideal $-1/2$)
 & $h_{\rm sat}\propto\rho_s^{-0.48}$ ($R^2{=}0.90$)
 & $a_0,\ \beta$
 & \textbf{weak} \\
Collapse \cite{ZambuziSR2026}
 & $\chi_{\rm eff}$ saturates; GOE $\langle r\rangle{\to}0.53$; decoh.\ $\propto\Delta x^{2}$
 & $\eta$ forced (percolation threshold)
 & $\eta^{\dagger}$ (geometric); $\hbar$
 & \textbf{convergence} \\
\bottomrule
\end{tabular}%
}
\end{table*}

\emph{Geometry.}---A localised source in the coarse-grained causal-density field, relaxed
under the Benincasa--Dowker action with no imposed profile, gives a Newtonian potential
$\theta\propto1/r$ (exterior exponent $-0.992$), a Poisson law for every source shape
(correlation $0.9998$), exact linear superposition (residual $1.8\times10^{-9}$), and the
Schwarzschild dilation $\sqrt{1-2M/r}$ to $0.2\%$ by chain counting---all with $G$ and the
metric absent from the generator \cite{MatterGravityPRD,TEICDoc1}. The substrate \emph{does} fix a
scale: the response $G_{\rm net}=A/M=0.93$ in lattice units, constant to five digits across
the source. But $G_{\rm net}\propto1/K$ with $K$ the action normalisation (the granularity
scale), so the absolute $G$ is the one external input.

\emph{Matter.}---On the $SU(2)$ ferromagnet, a topological soliton exists with baryon
number $B=1$ and measured spin-$\tfrac12$ fermionic statistics (exchange phase $=2\pi$
rotation $\in\mathbb{Z}_2$, to $5\times10^{-16}$); quantising its rotation with the full
inertia, one calibration (the $N$--$\Delta$ splitting) fixes the single dimensionless
coupling $e=5.39$, after which $\mu_p/\mu_n=-1.51$ (experiment $-1.46$) and the multiplet
degeneracies $[4,16,36]$ are parameter-free \cite{MatterGravityPRD,AdkinsNappiWitten1983}. On the
$SU(3)$ ferromagnet the same substrate confines ($V\sim\sigma r$, $\sigma>0$) and yields an
exactly degenerate meson octet with $c=\sqrt{2/3}$ \cite{SU3PRD}. The internal scales
($M_{\rm Sk}\approx146$--$207$, $\sigma=1.35\to0.33$, $\lambda_{\rm Sk}=a/\sqrt{120}$) are
in lattice units; the physical GeV scale $f_\pi$ is external.

\emph{Cosmology.}---This is the weakest link, and we mark it as such. The longitudinal
susceptibility of the ordered ferromagnet rises as $\chi_\parallel\sim h^{-0.37}$
(measured $-0.4\pm0.1$ against the ideal deep-MOND value $-1/2$), reproducing the form of
the MOND interpolation function and matching, in the deep-MOND limit, the response sector
of a superfluid dark-matter theory \cite{PaperIV}. Two honesty flags: finite-size effects
push the exponent toward zero, so the form is weak, not solid; and the transverse Goldstone
(magnon) is quadratic ($\propto X$), \emph{not} the fractional $X^{3/2}$ of the superfluid
phonon---so the loose identification ``magnon $=$ superfluid phonon'' is false, and the
equivalence is a deep-MOND \emph{limit} statement. The substrate again fixes an internal
scale---the response knee $h_{\rm sat}\propto\rho_s^{-0.48}$ ($R^2=0.90$), a genuine
network scale tracking the spin stiffness---while the absolute acceleration $a_0$ and the
coupling $\beta$ stay external.

\emph{Collapse (unplanned convergence).}---The most recent and most delicate domain. An
independent objective-collapse theory, Stochastic Rupture (SR) \cite{ZambuziSR2026},
defines a collapse order parameter $\chi_{\rm eff}$, a spectral noise of random-matrix
type, and a localisation rate. Fed only the \emph{bare} Poisson causal spectrum---no
ferromagnet, no threshold, no time direction inserted---the substrate reproduces the
\emph{forms}: $\chi_{\rm eff}$ saturates spontaneously to a dimension-fixed value
(reproducing SR's two-sector dichotomy, adjacency saturating while the Laplacian decays);
the causal operator's spectrum is a Dyson Brownian motion of GOE class
($\langle r\rangle\to0.53$); and the decoherence rate scales as $\Delta x^{2}$ (fitted
exponent $2.14$). We present this as \emph{convergence}, not derivation, for three honest
reasons: the $\Delta x^2$ and $\sigma_x^{-2}$ laws follow partly from a generic Laplacian
limit; complete positivity of the generator requires a Euclidean continuation (a choice);
and the threshold $\eta$ stays external---though in a sharper sense than the dimensional
constants. The substrate \emph{forces} $\eta$ to the percolation threshold of the causal
graph, a Molloy--Reed quantity (degree heterogeneity) raised by the graph's clustering,
which is non-generic (it sits $\sim8\sigma$ below the Erd\H{o}s--R\'enyi value in four
dimensions) and dimension-dependent; the form and the mechanism are derived, and only the
precise value's finite-size/clustering correction is missing. So $\eta$ is external as a
\emph{geometric} scale, not as a dimensional world constant. The forms emerge; the absolute
scale ($\hbar$) does not.

\emph{Why the pattern is structural---and the honest QCD analogy.}---In all four domains
the same boundary appears: the substrate derives dimensionless structure and internal
scales in lattice units, but never the lattice$\to$SI conversion. This is the position of
lattice gauge theory \cite{Wilson1974,Creutz1980}: dimensionless ratios and the running
are computed on the lattice, but a single physical input sets the unit. We draw the analogy
\emph{with} its decisive difference. In QCD the absolute scale $\Lambda_{\rm QCD}$ is
itself \emph{generated} from a dimensionless coupling by dimensional transmutation
\cite{GrossWilczek1973,Politzer1973,ColemanWeinberg1973}: a classically scale-free theory
manufactures a scale through the quantum running, and only the overall physical unit needs
calibration. The causal substrate has shown no such mechanism: its internal scales
($\rho_s$, $\sigma$, $h_{\rm sat}$, $M_{\rm Sk}$) ride on the fixed action normalisation
$K$, not on a running coupling, and attempts to fix the physical unit from within the
theory have failed by many orders. So the honest statement is not that a dimensionless
substrate \emph{cannot} generate a scale---QCD shows it can---but that \emph{this}
substrate has not exhibited the transmutation that would do so. The cheapest equilibrium
channel is now closed: we measured whether a correlation length diverges at the ordering
transition---the minimal prerequisite for the essential singularity of dimensional
transmutation---and it does not. The connected correlation per link is $1/z$ with no tail,
and the coordination $z$ grows with system size, so larger systems are \emph{more}
mean-field; the transition produces no divergent critical length to anchor a manufactured
scale (consistent with the volume-law, latent-heat-free character of the $SU(3)$ ordering
transition \cite{SU3PRD}). Transmutation in equilibrium is therefore excluded. A
non-equilibrium route has since been tested as well: a parametric drive holding the
geometry in a genuine non-equilibrium steady state (stationarity and broken detailed
balance both verified) produced an apparent scaling exponent that did not survive an
amplitude sweep at identical parameters---the driven substrate is mean-field too. And the
closure now has a structural form: for any connection rule on the Poisson sprinkling that
is a Lorentz-invariant function of pair or neighbourhood invariants, the expected
coordination diverges---by a Campbell--Mecke argument the marginal connection probability
depends only on the proper time, and the orbit of neighbours at fixed proper time is the
non-compact boost hyperboloid---while every tested escape that abandons the invariant
measure (sequential growth, dynamical triangulations in three and four dimensions) lands
in tree-like or Bethe-like graphs, mean-field by their loop structure. No invariant rule
can furnish the finite, loop-carrying connectivity that a divergent critical length
requires. We
add the methodological caveat already noted: the guard forbids inserting scales, so the
external-ness of $G$, $a_0$, $f_\pi$, $\eta$ is in part a property of the discipline, not a
proven impossibility.

\emph{Scales, with a mechanism.}---The scale sector---equilibrium and the driven
steady state tested---is now closed with a positive account of \emph{why}. A global rescaling of the action ($S\to KS$) does not fix
any mass ratio (the hierarchy scales as $\sqrt{K}$), so rescaling manufactures no scale. The
thresholds the network does produce---the collapse $\eta$, and by analogy $\hbar$---are
calculable graph properties (a percolation threshold set by degree heterogeneity and
clustering) but input-dependent (on dimension), hence not pinnable universal constants. And
the one known scale-manufacturing mechanism, dimensional transmutation, has nowhere to
anchor, the transition having no divergent critical length. The thesis ``forms derive,
scales do not'' is thus complete for the equilibrium sector, extends to the one
non-equilibrium regime tested, and now carries both a mechanism and a structural reason: the
substrate fixes forms and graph-thresholds, not absolute universal numbers, and the
table's geometric category ($\dagger$) names which kind of external each scale is.

\emph{Falsifiability.}---The claim is refuted by a single counterexample: a domain in which
the lattice$\to$SI conversion is fixed by the substrate, with no external calibration. This
has been searched and not found. A dedicated programme asked whether the galactic effective
theory's four numbers ($a_0$, $\beta$, the vector mass $m_A$, the slip) descend from the
network: the verdict was that each is calibrated, not derived---$a_0$'s internal correlate
$h_{\rm sat}$ even has the \emph{opposite} sign to the natural hypothesis, and $\beta$
misses the operating point by $48\times$ except at fine-tuned criticality. Direct attempts
to fix the absolute unit died decisively (the gravitational coupling inferred from a proton
radius missed by $\sim10^{8}$ orders; a $\sqrt{\rho}$ inheritance for $m_A$ failed at
$9.5\sigma$). Every such search has returned ``form derived, scale external,'' which is the
evidence for the pattern rather than against it. The same boundary holds from the other
side: the galactic effective theory exhibits it independently---its acceleration scale is an
external infrared boundary scale ($a_0\sim cH_0$) and its coupling is calibrated, not
derived---so the two theories prove ``forms derive, scales do not'' each in its own way. A
final test asked whether they are in fact \emph{one} theory at two resolutions, the
continuum vacuum being the coarse graining of the discrete network. They are not: under an
explicitly anti-isospectral gate (a shared functional form is never evidence of identity),
the identification is falsified in the vector sector---the effective theory needs a massive
Proca vector as its galactic motor, whereas the network forces a massless gauge sector by
the non-locality of its causal d'Alembertian---while the shared deep-MOND $X^{3/2}$ form is
isospectral, fixed by Milgrom's theorem for any microstructure. They are two distinct
theories that converge by symmetry, not one theory at two resolutions; a static coupling
between them is likewise infeasible on the present map.

\emph{Conclusion.}---A single Poisson causal substrate derives the dimensionless structure
of four unrelated domains---Newtonian gravity to $0.2\%$; a spin-$\tfrac12$ baryon and
$SU(3)$ confinement with parameter-free ratios; the deep-MOND form (weakly); and the
order parameter, random-matrix spectrum and $\Delta x^2$ localisation of objective collapse
(as convergence)---and in each it fixes internal scales in lattice units while leaving the
single lattice$\to$SI conversion, and thus every absolute scale $G$, $\hbar$, $a_0$,
$f_\pi$, $\eta$, external. The substrate is in the lattice-QCD position, with the honest
difference that no dimensional-transmutation mechanism generating the physical scale has
been identified here. The boundary is sharp, it is falsifiable, and it has survived every
search---a derivation of the form of physics from causal order, not of its scales.

\begin{acknowledgments}
This work uses only the author's own simulations except where it converges with the
independent Stochastic Rupture programme \cite{ZambuziSR2026}; the anti-circularity guard,
the validation gates, and the pre-registered decision criteria of the cited campaigns are
part of the build, and the weak (cosmology) and convergence-only (collapse) results are
flagged as such.
\end{acknowledgments}

\bibliographystyle{apsrev4-2}
\bibliography{TEIC_CONSOLIDATED}

\end{document}
